\section{Downstream Benchmarks and Settings}
\label{app:downstream}
\label{app:exp_details}

This appendix provides implementation details omitted from the main text for reproducibility.
Unless otherwise specified, the pretrained backbone is frozen and only PEFT parameters are updated.
For each benchmark, all compared methods use the same data split, target modules, optimizer, batch size, number of epochs, learning-rate schedule, and random seeds whenever applicable.
For \textbf{ProLoSA}, the sparse support is selected after compressed warmup using the accumulated sensitivity score defined in Appendix~\ref{sec:routed_sparse_adapter}.
Unless a benchmark specifies otherwise (Llama-3.1 and ViT below), the reported warmup ratio $\rho_{\mathrm{w}}$ refers to the score-free compressed-only phase $T_{\mathrm{pre}}=\lfloor\rho_{\mathrm{w}}T\rfloor$; the sparse branch is activated only after the additional fixed $T_{\mathrm{s}}=128$-step scoring window, i.e., at step $\lfloor\rho_{\mathrm{w}}T\rfloor+128$.
The selected support is fixed during the remaining training stage, and the sparse coefficients are initialized to zero.
All reported training time and memory statistics include method-specific stages, including the warmup stage of \textbf{ProLoSA}.


\subsection{Main Results}

\paragraph{Results on GLUE.}
Table~\ref{tab:glue} reports the results on GLUE \citep{wang2018glue}. 
ProLoSA improves the average score of Uni-LoRA \citep{li2025unilora} from 85.3 to 85.5 on RoBERTa-base and from 88.3 to 88.7 on RoBERTa-large with the same trainable budget, achieving the best average among compressed baselines.
The stronger relative performance of compressed methods on RoBERTa-large than on RoBERTa-base (where full LoRA retains an edge) is qualitatively compatible with the variance-reduction interpretation; however, changing the backbone alters not only the full-space dimension $D$ but also the representation quality, curvature, gradient noise, and target update simultaneously, so we do not treat this comparison as a controlled test of Eq.~\eqref{eq:energy_noise_risks}---the controlled mechanism checks are reported in Appendix~\ref{sec:real_model_validation}.

\begin{table*}[t]
\centering
\caption{Results with RoBERTa$_{\text{base}}$ and RoBERTa$_{\text{large}}$ on GLUE \citep{wang2018glue}: Matthews correlation (CoLA), Pearson correlation (STS-B), and accuracy (others); mean$_{\pm\mathrm{std}}$ over 5 seeds. Baseline results are from Uni-LoRA \citep{li2025unilora}. $^\dagger$For VB-LoRA we quote the trainable-parameter count reported by the original paper, whose counting convention (stored vs.\ trained parameters) differs from the other rows.}
\label{tab:glue}
\resizebox{0.9\linewidth}{!}{%
\begin{tabular}{clcccccccc}
\toprule
Backbone & Method & Trainable Params & SST-2 & MRPC & CoLA & QNLI & RTE & STS-B & Avg. \\
\midrule
\multirow{8}{*}{\rotatebox[origin=c]{90}{BASE}} 
& FT & 125M & 94.8 & 90.2 & 63.6 & 92.8 & 78.7 & 91.2 & 85.2 \\
& LoRA & 0.295M & $\textbf{95.1}_{\pm 0.2}$ & $89.7_{\pm 0.7}$ & $63.4_{\pm 1.2}$ & $\textbf{93.3}_{\pm 0.3}$ & $\textbf{86.6}_{\pm 0.7}$ & $\textbf{91.5}_{\pm 0.2}$ & \textbf{86.6} \\
\cmidrule{2-10}
& VeRA & 0.043M & $94.6_{\pm 0.1}$ & $89.5_{\pm 0.5}$ & $\textbf{65.6}_{\pm 0.8}$ & $91.8_{\pm 0.2}$ & $78.7_{\pm 0.7}$ & $90.7_{\pm 0.2}$ & 85.2 \\
& Tied-LoRA & 0.043M & $94.4_{\pm 0.5}$ & $88.5_{\pm 1.0}$ & $61.9_{\pm 1.6}$ & $92.0_{\pm 0.4}$ & $76.2_{\pm 1.3}$ & $89.8_{\pm 0.3}$ & 83.8 \\
& VB-LoRA & 0.075M$^\dagger$ & $94.4_{\pm 0.2}$ & $89.5_{\pm 0.5}$ & $63.3_{\pm 0.7}$ & $92.2_{\pm 0.2}$ & $82.3_{\pm 1.3}$ & $90.8_{\pm 0.1}$ & 85.4 \\
& FourierFT & 0.024M & $94.2_{\pm 0.3}$ & $\textbf{90.0}_{\pm 0.8}$ & $63.8_{\pm 1.6}$ & $92.2_{\pm 0.1}$ & $79.1_{\pm 0.5}$ & $90.8_{\pm 0.2}$ & 85.0 \\
& Uni-LoRA & 0.023M & $94.5_{\pm 0.2}$ & $89.7_{\pm 0.4}$ & $64.6_{\pm 1.0}$ & $92.3_{\pm 0.4}$ & $79.8_{\pm 1.1}$ & $90.9_{\pm 0.2}$ & 85.3 \\
& ProLoSA & 0.023M & $94.9_{\pm 0.3}$ & $89.5_{\pm 0.6}$ & $63.5_{\pm 0.2}$ & $92.5_{\pm 0.2}$ & $81.7_{\pm 0.6}$ & $90.8_{\pm 0.2}$ & 85.5 \\
\midrule
\multirow{8}{*}{\rotatebox[origin=c]{90}{LARGE}} 
& LoRA & 0.786M & $96.2_{\pm 0.5}$ & $90.2_{\pm 1.0}$ & $68.2_{\pm 1.9}$ & $\textbf{94.8}_{\pm 0.3}$ & $85.2_{\pm 1.1}$ & $\textbf{92.3}_{\pm 0.5}$ & 87.8 \\
\cmidrule{2-10}
& VeRA & 0.061M & $96.1_{\pm 0.1}$ & $90.9_{\pm 0.7}$ & $68.0_{\pm 0.8}$ & $94.4_{\pm 0.2}$ & $85.9_{\pm 0.7}$ & $91.7_{\pm 0.8}$ & 87.8 \\
& Tied-LoRA & 0.066M & $94.8_{\pm 0.6}$ & $89.7_{\pm 1.0}$ & $64.7_{\pm 1.2}$ & $94.1_{\pm 0.1}$ & $81.2_{\pm 0.1}$ & $90.8_{\pm 0.3}$ & 85.9 \\
& VB-LoRA & 0.162M$^\dagger$ & $96.1_{\pm 0.2}$ & $91.4_{\pm 0.6}$ & $68.3_{\pm 0.7}$ & $94.7_{\pm 0.5}$ & $86.6_{\pm 1.3}$ & $91.8_{\pm 0.1}$ & 88.2 \\
& LoRA-XS & 0.025M & $95.9_{\pm 0.3}$ & $90.7_{\pm 0.4}$ & $67.0_{\pm 1.2}$ & $93.9_{\pm 0.1}$ & $\textbf{88.1}_{\pm 0.3}$ & $92.0_{\pm 0.1}$ & 87.9 \\
& FourierFT & 0.048M & $96.0_{\pm 0.2}$ & $90.9_{\pm 0.3}$ & $67.1_{\pm 1.4}$ & $94.4_{\pm 0.4}$ & $87.4_{\pm 1.6}$ & $92.0_{\pm 0.4}$ & 88.0 \\
& Uni-LoRA & 0.023M & $96.3_{\pm 0.2}$ & $91.3_{\pm 0.6}$ & $68.5_{\pm 1.1}$ & $94.6_{\pm 0.4}$ & $86.6_{\pm 1.6}$ & $92.1_{\pm 0.1}$ & $88.3$ \\
& ProLoSA & 0.023M & $\textbf{96.5}_{\pm 0.2}$ & $\textbf{91.7}_{\pm 0.6}$ & $\textbf{70.3}_{\pm 0.6}$ & $94.6_{\pm 0.1}$ & $87.4_{\pm 0.6}$ & $91.9_{\pm 0.2}$ & \textbf{88.7} \\
\bottomrule
\end{tabular}%
}
\end{table*}


\paragraph{Results on mathematical reasoning.}
Table~\ref{tab:math} reports results for Gemma-7B \citep{team2024gemma} on GSM8K \citep{cobbe2021training} and MATH \citep{hendrycks2021measuring}, and for Llama-3.1-8B on GSM8K and MATH500. On Gemma-7B, ProLoSA shows its main advantage on the harder MATH benchmark, improving over Uni-LoRA from 28.94 to 29.90 under the same 0.52M budget, while remaining competitive on GSM8K, where Uni-LoRA is strongest among all methods (75.59) and ProLoSA's sparse branch does not pay off (74.69).
On Llama-3.1-8B, LoRA ($r=64$, 167.77M parameters) achieves 79.61/33.20 on GSM8K/MATH500. At a shared 1.05M budget, Uni-LoRA achieves 75.44/26.20 and ProLoSA achieves 75.06/26.60: ProLoSA scores 0.40 points higher on MATH500 and 0.38 points lower on GSM8K. These single-seed results show a tradeoff rather than a consistent improvement; both compressed methods remain below rank-64 LoRA. The displayed compressed configurations maximize the two-benchmark mean within each method at this budget among the available runs; this is a post-hoc test-set summary, with the full sweep reported in Appendix~\ref{app:llama31_math_sweep}.
We treat these results purely as a practical benchmark: because the model is fine-tuned once on MetaMathQA and then evaluated on both benchmarks (Appendix~\ref{app:exp_details}), the two evaluations share a single training update, so per-benchmark differences cannot be attributed to differences in task-specific update energy; a theoretical attribution would require training separately on each task's data and measuring the corresponding curvature-weighted off-subspace energies, which we leave to the controlled experiments of Appendix~\ref{sec:real_model_validation}.
\begin{table}[t]
\centering
\small
\begin{tabular}{lccc}
\toprule
\multicolumn{4}{c}{\textit{Gemma-7B}} \\
\midrule
Method & Trainable Params & GSM8K & MATH \\
\midrule
LoRA        & 200M    & {74.90} & \textbf{31.28} \\
VeRA        & 1.90M   & 74.98 & 28.84 \\
FourierFT   & 0.59M   & 72.97 & 25.14 \\
Uni-LoRA    & 0.52M   & \textbf{75.59} & 28.94 \\
{ProLoSA}   & 0.52M   & 74.69 & 29.90    \\
\midrule
\multicolumn{4}{c}{\textit{Llama-3.1-8B}} \\
\midrule
Method & Trainable Params & GSM8K & MATH500 \\
\midrule
LoRA ($r=64$) & 167.77M & \textbf{79.61} & \textbf{33.20} \\
Uni-LoRA & 1.05M & 75.44 & 26.20 \\
ProLoSA ($d:K=12:1$) & 1.05M & 75.06 & 26.60 \\
\bottomrule
\end{tabular}
\caption{Mathematical reasoning exact-match accuracy (\%). Gemma-7B uses full MATH; Llama-3.1-8B uses MATH500. Llama results use seed 42, with compressed configurations selected as described in the text; 1.05M denotes exactly 1,048,576 trainable parameters.}
\label{tab:math}
\end{table}




\paragraph{Results on instruction tuning.}
Table~\ref{tab:instruction} in Appendix~\ref{app:exp_details} summarizes instruction tuning under the QLoRA recipe \citep{dettmers2023qlora}. On multi-turn MT-Bench \citep{zheng2023judging}, ProLoSA outperforms both Uni-LoRA and standard LoRA \citep{hu2022lora} while using orders of magnitude fewer trainable parameters than the latter ($4.25$ vs.\ $3.23$ on LLaMA2-7B and $4.72$ vs.\ $4.13$ on LLaMA2-13B with $0.3\%$ of the LoRA budget); on single-turn scores it is competitive with Uni-LoRA but remains below LoRA on the 7B model ($5.36$ vs.\ $5.62$).


\subsection{Ablation Study}
\label{sec:ablation}

We ablate three design choices on CoLA and MRPC under matched trainable-parameter budgets: budget allocation between the two branches, sparse support selection, and warmup ratio; full results are reported in Appendix~\ref{app:ablation} (Table~\ref{tab:ablation_budget_full}).
Neither pure compression nor pure sparsity is optimal: hybrid allocations outperform the sparse-only variant by a large margin (e.g., $69.62$ vs.\ $59.99$ CoLA Matthews correlation at $d{+}K=23040$) and improve over compressed-only on at least one task.
Under the same $d=22320$ and $K=720$, warmup SNIP support selection outperforms both random and accumulated-gradient supports ($69.62$ vs.\ $66.58$ and $67.95$ on CoLA), where the accumulated-gradient variant scores coordinates by $\sum_t|g_i^{(t)}|$ without the $\theta$-weighting of Eq.~\eqref{eq:warmup_snip_score}.
Under the idealized isotropic matched-budget model, a random support is expected to provide no systematic gain over pure compression (Corollary~\ref{cor:random_support}); in the real network it performs slightly worse than compressed-only ($66.58$ vs.\ $68.99$), which may reflect anisotropic curvature, optimization and conditioning effects of the augmented parameterization, or finite-sample variability.
The consistent ordering SNIP $>$ random supports the central prediction that the gain comes from \emph{informed residual allocation} rather than from sparsity itself.
Finally, varying the score-free warmup phase $T_{\mathrm{pre}}=\lfloor\rho_{\mathrm{w}}T\rfloor$ that precedes the fixed $128$-step scoring window, a ratio of $10\%$ performs best, whereas $20\%$ degrades both tasks, likely because the sparse branch is activated too late and the joint-optimization phase becomes too short.


\subsection{GLUE Hyperparameters}

Table~\ref{tab:glue_hyperparams} reports the main hyperparameters for GLUE.
We use the same settings for all compared PEFT methods under the corresponding backbone unless otherwise stated.
For CoLA, STS-B, and the remaining tasks, we report Matthews correlation, Pearson correlation, and accuracy, respectively.
The median performance over five random seeds is reported.

\begin{table}[htbp]
\centering
\footnotesize
\setlength{\tabcolsep}{2pt}
\begin{tabular}{llcccccc}
\toprule
Backbone & Task & $d$ & $K$ & Epochs & BS & LR \\
\midrule
RoBERTa-base  & SST-2 & 22320 & 720 & 60  & 32 & 5e-4 \\
RoBERTa-base  & MRPC  & 22320 & 720 & 30  & 32 & 1e-3 \\
RoBERTa-base  & CoLA  & 14400 & 8640 & 80  & 32 & 2e-2 \\
RoBERTa-base  & QNLI  & 22320 & 720 & 25  & 32 & 1e-3 \\
RoBERTa-base  & RTE   & 22320 & 720 & 160 & 32 & 1e-2 \\
RoBERTa-base  & STS-B & 22320 & 720 & 80  & 32 & 2e-4 \\
RoBERTa-large & SST-2 & 22320 & 720 & 20  & 32 & 1e-3 \\
RoBERTa-large & MRPC  & 22320 & 720 & 40  & 32 & 2e-4 \\
RoBERTa-large & CoLA  & 14400 & 8640 & 40  & 32 & 5e-3 \\
RoBERTa-large & QNLI  & 22320 & 720 & 20  & 32 & 5e-4 \\
RoBERTa-large & RTE   & 22320 & 720 & 40  & 32 & 5e-3 \\
RoBERTa-large & STS-B & 22320 & 720 & 40  & 32 & 2e-4 \\
\midrule
Optimizer & \multicolumn{6}{c}{AdamW} \\
Scheduler & \multicolumn{6}{c}{Linear} \\
Seeds & \multicolumn{6}{c}{5} \\
ProLoSA warmup ratio & \multicolumn{6}{c}{10\%} \\
SNIP score-collection window & \multicolumn{6}{c}{128 steps}\\
\bottomrule
\end{tabular}
\caption{Core hyperparameter settings for GLUE experiments. BS denotes batch size. The $d$ and $K$ columns denote the compressed latent dimension and sparse residual budget used by ProLoSA.}
\label{tab:glue_hyperparams}
\end{table}

\subsection{Mathematical Reasoning Hyperparameters}

Table~\ref{tab:math_hyperparams} reports the Gemma-7B configuration for GSM8K and MATH.
Within each backbone, all methods use the same target modules, numerical precision, and decoding protocol; the Llama-3.1-8B settings and sweep are reported below.
Exact-match scores are computed with the same answer extraction rule for all methods.

\begin{table}[htbp]
\centering
\small
\setlength{\tabcolsep}{6pt}
\begin{tabular}{lcc}
\toprule
\textbf{Hyperparameter} & \textbf{GSM8K} & \textbf{MATH} \\
\midrule
\multicolumn{3}{c}{\textit{Shared Training Configuration}} \\
\midrule
Backbone & \multicolumn{2}{c}{Gemma-7B} \\
Train data & \multicolumn{2}{c}{MetaMathQA train} \\
Epochs & \multicolumn{2}{c}{2} \\
Train batch size & \multicolumn{2}{c}{1} \\
Gradient accumulation & \multicolumn{2}{c}{64} \\
Learning rate & \multicolumn{2}{c}{$2\times10^{-4}$} \\
Scheduler & \multicolumn{2}{c}{Cosine} \\
Max sequence length & \multicolumn{2}{c}{512} \\
ProLoSA $d$ & \multicolumn{2}{c}{349525} \\
ProLoSA $K$ & \multicolumn{2}{c}{174763} \\
ProLoSA warmup ratio & \multicolumn{2}{c}{10\%} \\
SNIP score-collection window & \multicolumn{2}{c}{128 steps}\\
Seeds & \multicolumn{2}{c}{42} \\
\midrule
\multicolumn{3}{c}{\textit{Task-Specific Evaluation Configuration}} \\
\midrule
Metric & \multicolumn{2}{c}{Exact match} \\
Eval batch size & 32 & 32 \\
Max new tokens & 1024 & 2048 \\
\bottomrule
\end{tabular}
\caption{Core training and evaluation settings for mathematical reasoning experiments. The model is fine-tuned once on MetaMathQA and evaluated across both benchmarks.}
\label{tab:math_hyperparams}
\end{table}


\paragraph{Llama-3.1-8B configuration and sweep.}
\label{app:llama31_math_sweep}
We fine-tune the Llama-3.1-8B base model (NousResearch mirror) on the first 100,000 MetaMathQA training examples for two epochs, using BF16 without weight quantization, a maximum training length of 512, batch size 1, gradient accumulation 64, cosine learning-rate decay, a 2\% scheduler warmup, and zero weight decay. All methods target the attention and MLP linear projections. LoRA uses $r=64$, $\alpha=64$, and learning rate $2\times10^{-4}$; Uni-LoRA and ProLoSA use $r=4$. Evaluation uses greedy decoding (temperature 0, top-$p=1$), batch size 32, and at most 1,024/2,048 generated tokens for GSM8K/MATH500. Each evaluated configuration uses seed 42 and the same fine-tuned checkpoint across the two benchmarks (1,319 and 500 examples, respectively); full MATH results are unavailable for this backbone.

In Table~\ref{tab:math}, Uni-LoRA uses $d=1{,}048{,}576$ and learning rate $2\times10^{-3}$. ProLoSA uses $d=967{,}916$, $K=80{,}660$ (approximately $12:1$), and a projected-vector learning rate of $1.6\times10^{-3}$. Its base learning rate is $2\times10^{-4}$, with sparse learning-rate multiplier 0.2, sparse warmup 128 steps, a one-step SNIP mask-collection window, and optimizer reset on mask activation. Its trainable budget counts both $d$ and $K$, including the sparse parameters activated after warmup.

Table~\ref{tab:llama31_math_sweep} lists all completed Llama-3.1-8B configurations found in the local evaluation results. At the shared 1,048,576-parameter budget, we display in the main table the configuration with the largest arithmetic mean of GSM8K and MATH500 for each compressed method. This descriptive selection uses the reported test scores, not an independent validation set; it does not establish statistical significance or a multi-seed advantage. The larger-budget Uni-LoRA runs and alternative ProLoSA ratios are retained below so that budget and hyperparameter effects remain visible.

\begin{table}[p]
\centering
\small
\setlength{\tabcolsep}{5pt}
\begin{tabular}{lrrrrr}
\toprule
Method & Params & $d:K$ & Learning rate & GSM8K & MATH500 \\
\midrule
LoRA ($r=4$) & 10,485,760 & -- & $2\times10^{-4}$ & 75.59 & 26.80 \\
LoRA ($r=64$) & 167,772,160 & -- & $2\times10^{-4}$ & 79.61 & 33.20 \\
\midrule
Uni-LoRA & 262,144 & -- & $2\times10^{-3}$ & 70.28 & 23.40 \\
Uni-LoRA & 524,288 & -- & $5\times10^{-4}$ & 68.39 & 25.40 \\
Uni-LoRA & 524,288 & -- & $1\times10^{-3}$ & 72.33 & 23.00 \\
Uni-LoRA & 524,288 & -- & $2\times10^{-3}$ & 73.84 & 24.80 \\
Uni-LoRA & 524,288 & -- & $4\times10^{-3}$ & 74.15 & 26.60 \\
Uni-LoRA & 524,288 & -- & $8\times10^{-3}$ & 74.75 & 24.20 \\
Uni-LoRA & 524,288 & -- & $1\times10^{-2}$ & 72.40 & 21.00 \\
Uni-LoRA & 1,048,576 & -- & $5\times10^{-4}$ & 72.25 & 25.00 \\
Uni-LoRA & 1,048,576 & -- & $1\times10^{-3}$ & 74.22 & 26.20 \\
Uni-LoRA & 1,048,576 & -- & $2\times10^{-3}$ & 75.44 & 26.20 \\
Uni-LoRA & 1,048,576 & -- & $4\times10^{-3}$ & 74.07 & 25.60 \\
Uni-LoRA & 2,097,152 & -- & $5\times10^{-4}$ & 74.45 & 26.00 \\
Uni-LoRA & 2,097,152 & -- & $1\times10^{-3}$ & 75.97 & 25.60 \\
Uni-LoRA & 2,097,152 & -- & $2\times10^{-3}$ & 76.95 & 26.00 \\
Uni-LoRA & 2,097,152 & -- & $4\times10^{-3}$ & 74.30 & 25.80 \\
Uni-LoRA & 4,194,304 & -- & $2\times10^{-3}$ & 76.12 & 26.20 \\
\midrule
ProLoSA & 524,288 & 2:1 & $8\times10^{-4}$ & 69.75 & 24.80 \\
ProLoSA & 524,288 & 4:1 & $8\times10^{-4}$ & 69.07 & 21.40 \\
ProLoSA & 524,288 & 8:1 & $8\times10^{-4}$ & 69.83 & 21.00 \\
ProLoSA & 524,288 & 12:1 & $4\times10^{-4}$ & 68.23 & 21.00 \\
ProLoSA & 524,288 & 12:1 & $8\times10^{-4}$ & 71.27 & 21.00 \\
ProLoSA & 524,288 & 12:1 & $1.6\times10^{-3}$ & 73.31 & 20.60 \\
ProLoSA & 524,288 & 12:1 & $3.2\times10^{-3}$ & 73.92 & 23.80 \\
ProLoSA & 1,048,576 & 2:1 & $1.6\times10^{-3}$ & 73.77 & 24.80 \\
ProLoSA & 1,048,576 & 2:1 & $3.2\times10^{-3}$ & 74.98 & 24.60 \\
ProLoSA & 1,048,576 & 4:1 & $1.6\times10^{-3}$ & 75.66 & 24.80 \\
ProLoSA & 1,048,576 & 4:1 & $3.2\times10^{-3}$ & 74.60 & 25.00 \\
ProLoSA & 1,048,576 & 8:1 & $1.6\times10^{-3}$ & 74.07 & 26.20 \\
ProLoSA & 1,048,576 & 8:1 & $3.2\times10^{-3}$ & 74.37 & 25.60 \\
ProLoSA & 1,048,576 & 12:1 & $1.6\times10^{-3}$ & 75.06 & 26.60 \\
ProLoSA & 1,048,576 & 12:1 & $3.2\times10^{-3}$ & 74.53 & 21.60 \\
\bottomrule
\end{tabular}
\caption{Complete available Llama-3.1-8B sweep (exact-match accuracy in \%, seed 42). The learning-rate column gives the projected-vector rate for ProLoSA, whose base rate is fixed at $2\times10^{-4}$, and the optimizer rate for LoRA/Uni-LoRA. Ratios $d:K$ are approximate due to integer budgets. Rank-4 LoRA is included only as a supplementary reference; the main baseline uses rank 64. Each row reports both benchmarks from one configuration.}
\label{tab:llama31_math_sweep}
\end{table}

\subsection{Instruction-Tuning Hyperparameters}

Table~\ref{tab:instruction} reports the full instruction-tuning results referenced in the main text, and Table~\ref{tab:instruction_hyperparams} reports the core configuration.
All methods are trained on the same instruction data and evaluated with the same MT-Bench protocol.
For QLoRA-style experiments, the backbone is quantized and frozen, and only PEFT parameters are updated.

\begin{table}[htbp]
\centering
\small
\begin{tabular}{lccc}
\toprule
\textbf{Method} & \textbf{\# Params} & \textbf{Score$_1$} & \textbf{Score$_2$} \\
\midrule
w/o FT       & -       & 1.31 / 1.46 & 1.11 / 1.06 \\
LoRA$^{\#}$  & 159.9M / 250.3M & \textbf{5.62} / 6.20 & 3.23 / 4.13 \\
Uni-LoRA     & 0.52M / 1.0M    & 5.34 / 6.34 & 3.56 / 4.43 \\
ProLoSA      & 0.52M / 1.0M    & 5.36 / \textbf{6.57} & \textbf{4.25} / \textbf{4.72} \\
\bottomrule
\end{tabular}
\caption{
Instruction-tuning results under the QLoRA training recipe \citep{dettmers2023qlora}. 
Each entry reports LLaMA2-7B / LLaMA2-13B results. 
Score$_1$ and Score$_2$ denote MT-Bench single-turn and multi-turn scores, respectively.
}
\label{tab:instruction}
\end{table}

\begin{table}[t]
\centering
\small
\renewcommand{\arraystretch}{0.9} % 稍微压缩行距，节省纵向空间
\setlength{\tabcolsep}{4pt}
\begin{tabular}{lcc}
\toprule
Hyperparameter & Shared / LLaMA2-7B & LLaMA2-13B \\
\midrule
Training data & \multicolumn{2}{c}{Cleaned Alpaca} \\
Evaluation & \multicolumn{2}{c}{MT-Bench} \\
Quantization & \multicolumn{2}{c}{4-bit NF4 + double quantization} \\
Epochs & \multicolumn{2}{c}{1} \\
Effective batch size & \multicolumn{2}{c}{16} \\
Per-device batch size & 4 & 2 \\
Gradient accumulation & 4 & 8 \\
Scheduler & \multicolumn{2}{c}{Linear} \\
Max sequence length & \multicolumn{2}{c}{16 / 512} \\
Decoding strategy & \multicolumn{2}{c}{Category-specific sampling/greedy} \\
Temperature & \multicolumn{2}{c}{0.7 / 0.1 / 0 by category} \\
Max new tokens & \multicolumn{2}{c}{1024} \\
Evaluator / judge & \multicolumn{2}{c}{GPT-4 single} \\
LoRA params & 159.9M & 250.3M \\
Uni-LoRA params & 0.52M & 1.0M \\
ProLoSA params & 0.52M & 1.0M \\
ProLoSA $d$ & 507904 & 983040 \\
ProLoSA $K$ & 16384 & 65536 \\
ProLoSA warmup ratio & \multicolumn{2}{c}{10\%} \\
SNIP score-collection window & \multicolumn{2}{c}{128 steps}\\
Seeds & \multicolumn{2}{c}{0} \\
\bottomrule
\end{tabular}
\caption{Core training and evaluation settings for instruction-tuning experiments. Score$_1$ and Score$_2$ denote MT-Bench single-turn and multi-turn scores, respectively.}
\label{tab:instruction_hyperparams}
\end{table}

\subsection{Ablation Details}
\label{app:ablation}

Table~\ref{tab:ablation_budget_full} reports the full budget-allocation sweep between the projected branch and the sparse residual branch.
All variants use the same total trainable-parameter budget $B=d+K=23040$.
The compressed-only variant corresponds to $K=0$, the sparse-only variant corresponds to $d=0$, and intermediate configurations correspond to \textbf{ProLoSA} with both branches active.
Since $B$ is fixed, increasing the sparse budget $K$ necessarily reduces the compressed latent dimension $d$.

\begin{table}[t]
\centering
\small
\renewcommand{\arraystretch}{0.9} % 纵向稍微压缩行距，节省高度
\setlength{\tabcolsep}{2.5pt}     % 缩小列与列之间的空白间距，避免缩小字号
\begin{tabular}{@{}llcccc@{}}      % @{} 去除表格最左和最右的额外空白占位
\toprule
Study & Variant & $d$ & $K$ & CoLA & MRPC \\
\midrule
\multirow{9}{*}{Budget}
& Compressed-only & 23040 & 0     & $68.99_{\pm 0.58}$ & $90.93_{\pm 0.00}$ \\
& ProLoSA & 22680 & 360   & $68.29_{\pm 1.40}$ & $90.11_{\pm 0.50}$ \\
& ProLoSA & 22320 & 720   & $69.62_{\pm 1.84}$ & $\mathbf{91.69}_{\pm 0.40}$ \\
& ProLoSA & 21600 & 1440  & $67.21_{\pm 0.84}$ & $90.87_{\pm 0.46}$ \\
& ProLoSA & 20160 & 2880  & $68.72_{\pm 0.45}$ & $83.74_{\pm 9.65}$ \\
& ProLoSA & 18720 & 4320  & $67.51_{\pm 0.51}$ & $90.03_{\pm 0.42}$ \\
& ProLoSA & 14400 & 8640  & $\mathbf{70.27}_{\pm 0.60}$ & $90.77_{\pm 0.42}$ \\
& ProLoSA & 5760  & 17280 & $68.75_{\pm 1.27}$ & $89.71_{\pm 0.20}$ \\
& Sparse-only & 0 & 23040 & $59.99_{\pm 1.40}$ & $84.89_{\pm 0.83}$ \\
\midrule
\multirow{3}{*}{Support}
& Random & 22320 & 720 & $66.58_{\pm 0.63}$ & $90.28_{\pm 0.61}$ \\
& Accum. gradient & 22320 & 720 & $67.95_{\pm 0.59}$ & $90.03_{\pm 0.31}$ \\
& Warmup SNIP & 22320 & 720 & $\mathbf{69.62}_{\pm 1.84}$ & $\mathbf{91.69}_{\pm 0.40}$ \\
\midrule
\multirow{4}{*}{Warmup}
& $0\%$ & 22320 & 720 & $67.32_{\pm 0.47}$ & $90.36_{\pm 0.20}$ \\
& $5\%$ & 22320 & 720 & $69.05_{\pm 0.51}$ & $91.12_{\pm 0.70}$ \\
& $10\%$ & 22320 & 720 & $\mathbf{69.62}_{\pm 1.84}$ & $\mathbf{91.69}_{\pm 0.40}$ \\
& $20\%$ & 22320 & 720 & $67.49_{\pm 0.96}$ & $90.60_{\pm 0.81}$ \\
\bottomrule
\end{tabular}
\caption{
Ablation details on budget allocation, sparse support selection, and warmup length.
CoLA reports Matthews correlation and MRPC reports accuracy.
}
\label{tab:ablation_budget_full}
\end{table}

\clearpage
\subsection{Vision Transformer Experiments}
\label{app:vit}

\paragraph{Relation to the original vision experiments.}
Uni-LoRA \citep{li2025unilora} already evaluates ViT-Base and ViT-Large on eight image-classification datasets, with five runs per configuration (its Section~4.4 and Table~5). We extend this line of evaluation to ProLoSA using the completed ViT-B/16 runs. Our data subsets, optimization settings, and seed coverage differ from that study, so all LoRA and Uni-LoRA scores in Table~\ref{tab:vit} come from the same local experiment suite as ProLoSA. They are not copied from the original paper. No completed ProLoSA ViT-Large results are available in this suite.

\paragraph{Model and data.}
We use the ImageNet-21k-pretrained ViT-B/16 checkpoint (\path{google/vit-base-patch16-224-in21k}) and apply rank-4 adapters to query and value projections in all 12 Transformer blocks. All other backbone parameters are frozen, and a randomly initialized classification head is trained for every method. The three limited-data settings use 1,000 training images from CIFAR-100, Food-101, and DTD, respectively. Sampling is stratified by class with subset seed 42, giving 10 images per class for CIFAR-100 and approximately balanced allocations for Food-101 and DTD. CIFAR-100 additionally uses 5,000 images and the full 50,000-image training split. These are custom stratified subsets, not the standard VTAB-1k protocol. Evaluation uses the complete CIFAR-100 test split (10,000 images), Food-101 validation split (25,250), and DTD test split (1,880), respectively. Training applies random resized cropping to $224\times224$, horizontal flipping, and pretrained-model normalization; evaluation uses resizing, center cropping, and the same normalization.

\paragraph{Budgets and parameter accounting.}
The rank-4 LoRA factor space contains $D=12\times2\times(768\times4+4\times768)=147{,}456$ parameters. Uni-LoRA uses $d=24{,}600$, while ProLoSA divides the same budget into $d+K=24{,}600$ (Table~\ref{tab:vit_allocations}). The shared head adds $769C$ parameters for $C$ classes: 76,900 for CIFAR-100, 77,669 for Food-101, and 36,143 for DTD. Thus the total active parameter counts for either compressed method are 101,500, 102,269, and 60,743, respectively; LoRA totals are 224,356, 225,125, and 183,599. Adapter-only compression is approximately $6\times$; this factor does not apply to the totals including the head. ProLoSA counts the $K$ selected sparse coefficients even though they are frozen during warmup. The raw result files count parameters before activation and therefore contain only $d$ in their ProLoSA adapter field; we recover $d+K$ from the logged configuration and verify $K$ against the activation log. These counts describe active degrees of freedom, not allocated tensor storage or optimizer memory.

\begin{table}[t]
\centering\small
\begin{tabular}{lrr@{\qquad}lrr}
\toprule
\multicolumn{3}{c}{Main budget: $d+K=24{,}600$} & \multicolumn{3}{c}{Supplementary budget: $d+K=72{,}000$} \\
\cmidrule(lr){1-3}\cmidrule(lr){4-6}
Method & $d$ & $K$ & Method & $d$ & $K$ \\
\midrule
Uni-LoRA & 24,600 & 0 & Uni-LoRA & 72,000 & 0 \\
ProLoSA ($2:1$) & 16,400 & 8,200 & ProLoSA ($2:1$) & 48,000 & 24,000 \\
ProLoSA ($4:1$) & 19,680 & 4,920 & ProLoSA ($4:1$) & 57,600 & 14,400 \\
ProLoSA ($8:1$) & 21,867 & 2,733 & ProLoSA ($8:1$) & 64,000 & 8,000 \\
ProLoSA ($12:1$) & 22,708 & 1,892 & ProLoSA ($12:1$) & 66,462 & 5,538 \\
\bottomrule
\end{tabular}
\caption{Exact compressed-adapter budgets for ViT-B/16. Ratios are approximate when integer rounding is required. The task-specific classification head is additional and shared across methods.}
\label{tab:vit_allocations}
\end{table}

\paragraph{Optimization and support selection.}
All runs use AdamW with $(\beta_1,\beta_2)=(0.9,0.999)$, $\epsilon=10^{-8}$, zero weight decay, BF16, batch size 64, and no gradient accumulation. The 1,000-image settings train for 20 epochs (320 optimizer steps); CIFAR-100 with 5,000 and 50,000 images trains for 10 epochs (790 and 7,820 steps). LoRA uses $\alpha=r=4$ and adapter learning rate $10^{-3}$; Uni-LoRA and ProLoSA use a latent-vector learning rate of $4\times10^{-3}$. The head learning rate is $3\times10^{-3}$ for every method. Adapter dropout is zero. The latent/head learning rates follow linear decay with 10\% scheduler warmup. ProLoSA uses latent initialization bound 0.02, a score-free warmup of 32 steps for the 1,000-image settings or 128 steps otherwise, and a one-step SNIP scoring window. Its support is activated after step 33 or 129, respectively, with zero-initialized sparse coefficients and cleared optimizer state. The sparse learning rate is $0.2$ times the latent rate ($8\times10^{-4}$), with linear decay after activation. This support-selection schedule differs from the 128-step scoring window used in the earlier language experiments.

\paragraph{Reporting and limitations.}
We report \emph{final-model} top-1 accuracy, using the \texttt{final\_accuracy} field rather than \texttt{best\_accuracy}, which is the maximum over repeated evaluations on the same held-out split. The full method-by-setting comparison is available only for training seed 42, and Table~\ref{tab:vit} retains all four ProLoSA ratios. Only LoRA on CIFAR-100 with 1,000 images has additional completed seeds: its final accuracies for seeds 42/43/44 are 86.91/86.57/86.44. We therefore do not attach multi-seed error bars to the compressed-method comparison. No independent validation split was used to choose between the available budget/learning-rate recipes, so the highlighted recipe is an exploratory result. Full-data CIFAR-100 improves by 0.23 points over Uni-LoRA at $d:K=4:1$, but the four ProLoSA ratios lose 0.09--0.25 points at 1,000 images and 0.19--0.42 points at 5,000 images. The small Food-101 and DTD differences and the lack of repeated ProLoSA seeds preclude a consistent-advantage claim. The changing epoch count and number of updates across data sizes also prevent a controlled test of the sample-size law.

\paragraph{Supplementary 72,000-parameter recipe.}
Table~\ref{tab:vit_72k} retains the second complete single-seed sweep. It uses $d=72{,}000$ for Uni-LoRA and $d+K=72{,}000$ for ProLoSA, with latent learning rate $10^{-2}$, sparse learning rate $2\times10^{-3}$, and head learning rate $10^{-2}$ on CIFAR-100 or $5\times10^{-2}$ on Food-101 and DTD. LoRA's adapter learning rate remains $10^{-3}$; data, epochs, and support-selection timing are unchanged. Since the learning rates change together with the compressed budget (including the head rate for LoRA), differences from Table~\ref{tab:vit} do not isolate a budget effect. At this recipe, ProLoSA's $2:1$ allocation exceeds Uni-LoRA on the four limited-data settings but loses on full-data CIFAR-100; LoRA remains stronger than this allocation on Food-101, DTD, and the larger CIFAR-100 subsets. This recipe dependence reinforces the need for independent tuning and repeated seeds.

\begin{table}[t]
\centering\small
\setlength{\tabcolsep}{5pt}
% Generated from final_accuracy, seed 42; do not edit.
\begin{tabular}{lrrrrr}
\toprule
 & \multicolumn{3}{c}{1,000 training images} & \multicolumn{2}{c}{CIFAR-100} \\
\cmidrule(lr){2-4}\cmidrule(lr){5-6}
Method & CIFAR-100 & Food-101 & DTD & 5,000 & Full \\
\midrule
LoRA ($r=4$) & 85.22 & 73.52 & 68.83 & 89.32 & 91.78 \\
Uni-LoRA & 85.10 & 72.68 & 66.70 & 88.24 & 91.51 \\
ProLoSA ($2:1$) & 85.38 & 72.77 & 67.93 & 88.96 & 91.43 \\
ProLoSA ($4:1$) & 85.08 & 71.95 & 68.19 & 88.55 & 91.80 \\
ProLoSA ($8:1$) & 84.79 & 72.48 & 68.09 & 88.60 & 91.57 \\
ProLoSA ($12:1$) & 84.98 & 72.37 & 67.93 & 88.72 & 91.76 \\
\bottomrule
\end{tabular}

\caption{Supplementary ViT-B/16 recipe: final-model accuracy (\%, seed 42), with 72,000 active adapter parameters for both compressed methods and 147,456 for LoRA, plus the shared classification head. All available ProLoSA ratios are retained. Learning rates differ from the main vision table as specified in the text.}
\label{tab:vit_72k}
\end{table}

